\chapter{Objetivos}\label{chap:objetivos}

En este capítulo se concretan los objetivos que orientan el Trabajo de Fin de Grado. A
partir de la motivación y el contexto expuestos en el Capítulo~\ref{chap:introduccion},
se formula primero el objetivo general y, a continuación, se desglosa en una serie de
objetivos específicos que estructuran el desarrollo experimental y guían la evaluación
de los \textit{frameworks} comparados.

\section{Objetivo general}\label{sec:objetivo-general}

El objetivo general de este Trabajo de Fin de Grado es comparar de forma experimental
varios \textit{frameworks} de aprendizaje federado. Para ello se ejecuta un mismo
escenario de entrenamiento en todos y se observa tanto el rendimiento del modelo como el
consumo de recursos y la facilidad de uso.

\section{Objetivos específicos}\label{sec:objetivos-especificos}

Para alcanzar este objetivo general se plantean los siguientes objetivos específicos:

\begin{enumerate}
  \item \textbf{OE1:} Estudiar los fundamentos del aprendizaje federado, su
        funcionamiento, sus diferencias respecto al entrenamiento centralizado y sus
        principales retos técnicos.
  \item \textbf{OE2:} Analizar la necesidad de utilizar \textit{frameworks}
        especializados para implementar sistemas federados e identificar los aspectos
        que influyen en su elección.
  \item \textbf{OE3:} Diseñar un escenario experimental común que permita comparar
        herramientas bajo condiciones equivalentes, manteniendo constantes el conjunto
        de datos, el modelo, los hiperparámetros y el entorno de ejecución en la medida
        de lo posible.
  \item \textbf{OE4:} Implementar el experimento federado en los \textit{frameworks}
        seleccionados, adaptando el flujo de entrenamiento a la arquitectura propia de
        cada herramienta.
  \item \textbf{OE5:} Recoger métricas de aprendizaje, rendimiento y consumo de
        recursos: precisión, pérdida, tiempo de ejecución, uso de GPU, memoria empleada,
        potencia consumida y energía estimada.
  \item \textbf{OE6:} Evaluar la experiencia práctica de uso de cada \textit{framework}:
        instalación, configuración, documentación, claridad de los registros, facilidad
        de depuración y flexibilidad.
  \item \textbf{OE7:} Comparar los resultados de forma crítica y señalar ventajas,
        limitaciones y escenarios en los que cada herramienta puede resultar más
        adecuada.
\end{enumerate}
